\chapter{Grundlagen}
\label{chap:grundlagen}

\section{Scheduling auf identischen Maschinen}
\label{sec:problem}

\begin{definition}[Instanz und Makespan]
\label{def:instanz}
Eine \emph{Instanz} $\Gamma$ besteht aus $n$ Aufgaben mit Größen $p_1, \ldots, p_n > 0$.
Eine \emph{Zuordnung} von $\Gamma$ auf $m$ Maschinen ist ein Tupel $P = (P_1, \ldots, P_m)$
paarweise disjunkter Mengen mit
\[
  P_1 \uplus P_2 \uplus \cdots \uplus P_m \;=\; \{1, \ldots, n\} .
\]
Die \emph{Last} einer Maschine $i$ ist $l_i = \sum_{j \in P_i} p_j$, der \emph{Makespan} der
Zuordnung ist $C_{\max}(P) = \max_{1 \le i \le m} l_i$.
Der kleinste Makespan einer Zuordnung von $\Gamma$ auf $m$ Maschinen wird mit
$C^{*}_{\max}(\Gamma, m)$ bezeichnet.
\end{definition}

Zu einer Eingabe des Scheduling-Problems gehören damit stets beide Angaben: die Aufgabenliste
$\Gamma$ und die Maschinenzahl $m$.
Sie werden hier getrennt geführt und nicht zu einem Objekt zusammengefasst, weil dieselbe
Aufgabenliste in Abschnitt~\ref{sec:binpacking} noch einmal auftritt - dort als
Bin-Packing-Instanz, zu der keine Maschinenzahl gehört.
Jede Aussage nennt die Maschinenzahl deshalb ausdrücklich, sobald sie gebraucht wird.

In der Drei-Feld-Notation $\alpha \mid \beta \mid \gamma$ von \textcite{graham1979} heißt dieses
Problem $P\,\|\,C_{\max}$: Das Feld $\alpha = P$ steht für identische parallele Maschinen, das
leere Feld $\beta$ dafür, dass an die Aufgaben keine weiteren Einschränkungen gestellt werden, und
$\gamma = C_{\max}$ für den Makespan als Zielgröße.
Das $P$ dieser Notation bezeichnet die Maschinenumgebung und ist von der Zuordnung $P$ aus
Definition~\ref{def:instanz} zu unterscheiden.
Die Aufgaben werden dabei weder unterbrochen noch auf mehrere Maschinen aufgeteilt.

Das zugehörige Entscheidungsproblem - gibt es zu einer gegebenen Schranke $T > 0$ eine Zuordnung
mit Makespan höchstens $T$? - ist NP-vollständig \cite{garey_johnson1979}.
Sofern $\mathrm{P} \ne \mathrm{NP}$ gilt, berechnet daher kein Algorithmus mit polynomieller
Laufzeit den Wert $C^{*}_{\max}(\Gamma, m)$.
Bekannt sind aber die beiden folgenden Schranken, auf die sich die Kapitel~\ref{kap:alternative}
und~\ref{chap:evaluation} wiederholt stützen.
Dabei bezeichnet $p_{\max} := \max_j p_j$ durchgehend die größte Aufgabengröße.

\begin{lemma}[Untere Schranken]
\label{lem:untere-schranken}
Für jede Instanz $\Gamma$ und jede Maschinenzahl $m \ge 1$ gilt
\[
  C^{*}_{\max}(\Gamma, m) \;\ge\; \max\Bigl( \tfrac{1}{m}\textstyle\sum_{j} p_j,\; p_{\max} \Bigr).
\]
\end{lemma}

\begin{proof}
Jede Zuordnung verteilt die Gesamtlast $\sum_j p_j$ auf $m$ Maschinen, sodass mindestens eine
Maschine eine Last von mindestens $\frac{1}{m}\sum_j p_j$ trägt.
Ferner liegt die größte Aufgabe in genau einem $P_i$, dessen Last $l_i$ damit mindestens
$p_{\max}$ beträgt.
Beides gilt für jede Zuordnung, insbesondere für eine mit minimalem Makespan.
\end{proof}

\begin{definition}[Approximationsgüte]
\label{def:guete}
Sei $A$ ein Algorithmus, der zu einer Instanz $\Gamma$ und einer Maschinenzahl $m$ eine Zuordnung
berechnet, und sei $C^{A}_{\max}(\Gamma, m)$ deren Makespan.
Die \emph{Güte} von $A$ auf $m$ Maschinen ist
\[
  R_m(A) \;=\; \sup\Bigl\{\, C^{A}_{\max}(\Gamma, m) \,/\, C^{*}_{\max}(\Gamma, m)
  \;\Bigm|\; \Gamma \text{ Instanz} \,\Bigr\}.
\]
\end{definition}

Steht die Maschinenzahl im Kontext fest, so wird das zweite Argument weggelassen und kurz
$C^{*}_{\max}(\Gamma)$ beziehungsweise $C^{A}_{\max}(\Gamma)$ geschrieben.

\section{Bin Packing und FFD}
\label{sec:binpacking}

\begin{definition}[Bin Packing]
\label{def:binpacking}
Zu einer Instanz $\Gamma$ und einer Kapazität $C > 0$ fragt das Bin-Packing-Problem nach der
kleinsten Anzahl von Bins der Kapazität $C$, auf die sich die Aufgaben so verteilen lassen, dass
kein Bin eine Last über $C$ trägt.
Diese Anzahl wird mit $\mathrm{OPT}[\Gamma, C]$ bezeichnet.
\end{definition}

Bin Packing kommt ohne Maschinenzahl aus: Gesucht ist die Anzahl der Bins, nicht eine Verteilung
auf eine vorgegebene Anzahl von Maschinen.
Beide Probleme arbeiten deshalb auf derselben Instanz $\Gamma$; über die Maschinenzahl sind sie
zueinander dual.

\begin{lemma}[Dualität]
\label{lem:dualitaet}
Sei $\Gamma$ eine Instanz, $m \ge 1$ eine Maschinenzahl und $C, C' > 0$ Kapazitäten.
Dann gilt $C^{*}_{\max}(\Gamma, m) \le C$ genau dann, wenn $\mathrm{OPT}[\Gamma, C] \le m$.
Ferner ist $C \mapsto \mathrm{OPT}[\Gamma, C]$ monoton fallend, das heißt aus $C \le C'$ folgt
$\mathrm{OPT}[\Gamma, C'] \le \mathrm{OPT}[\Gamma, C]$.
\end{lemma}

\begin{proof}
Eine Zuordnung mit Makespan höchstens $C$ verteilt die Aufgaben auf $m$ Mengen mit Last jeweils
höchstens $C$, ist also eine Packung in höchstens $m$ Bins der Kapazität $C$; leere Maschinen
entsprechen ungenutzten Bins.
Umgekehrt liefert eine Packung in höchstens $m$ solche Bins eine Zuordnung auf $m$ Maschinen mit
Makespan höchstens $C$.
Die Monotonie folgt, da jede Packung mit Kapazität $C$ auch für jede Kapazität $C' \ge C$ zulässig
ist.
\end{proof}

\emph{First Fit} legt jede Aufgabe in einer gegebenen Reihenfolge in den ersten Bin, in den sie
passt, und eröffnet einen neuen Bin, falls keiner passt.
\emph{First Fit Decreasing} (FFD) wendet First Fit auf die absteigend sortierte
Aufgabenliste an.
Mit $\mathrm{FFD}[\Gamma, C]$ wird die Anzahl der dabei erzeugten Bins bezeichnet.
Es gilt
\begin{equation}
  \mathrm{FFD}[\Gamma, C] \;\le\; \tfrac{11}{9}\,\mathrm{OPT}[\Gamma, C] + \tfrac{6}{9},
  \label{eq:ffd-guete}
\end{equation}
und diese Schranke ist scharf \cite{dosa2007}.

\section{MULTIFIT}
\label{sec:multifit}

MULTIFIT \cite{coffman1978} löst $P\,\|\,C_{\max}$ nicht direkt, sondern führt eine Binärsuche über
die Kapazität $C$ aus und ruft in jeder Iteration FFD auf.
Ausgehend von
\begin{equation}
  c_{\mathrm{low}}^{(0)} = \max\Bigl( \tfrac{1}{m}\textstyle\sum_j p_j,\; p_{\max} \Bigr),
  \qquad
  c_{\mathrm{high}}^{(0)} = \tfrac{1}{m}\textstyle\sum_j p_j + p_{\max}
  \label{eq:startintervall}
\end{equation}
wird $k$-mal die Intervallmitte $C = \frac{1}{2}(c_{\mathrm{low}} + c_{\mathrm{high}})$ getestet:
Gilt $\mathrm{FFD}[\Gamma, C] \le m$, so wird $c_{\mathrm{high}} := C$ gesetzt und die erzeugte
Packung gespeichert, andernfalls $c_{\mathrm{low}} := C$.
Zurückgegeben wird die zuletzt gespeicherte Packung; ihr Makespan wird mit
$C^{\mathrm{MF}}_{\max}(\Gamma)$ bezeichnet, die Variante mit $k$ Halbierungsschritten mit
$\mathrm{MF}(k)$.
Ein Aufruf mit $c_{\mathrm{high}}^{(0)}$ vor der Schleife stellt sicher, dass stets eine
gespeicherte Packung vorliegt; er ist der Rückfall für den Fall, dass alle $k$ Tests scheitern,
und lässt sich ebenso ans Ende verschieben.
Als obere Grenze wählen \textcite{coffman1978} ursprünglich
$\max\bigl(\tfrac{2}{m}\sum_j p_j,\, p_{\max}\bigr)$ \cite[S.~37]{friesen_langston1986}; diese
Arbeit verwendet durchgehend die Grenze aus~\eqref{eq:startintervall}.

Dass beide Grenzen aus \eqref{eq:startintervall} zulässig gewählt sind, hält das folgende Lemma
fest: Die untere Grenze überschreitet den gesuchten Wert nicht, und an der oberen gelingt der
Test.

\begin{lemma}[Startintervall]
\label{lem:startintervall}
Für jede Instanz $\Gamma$ und jede Maschinenzahl $m \ge 1$ gilt
\begin{enumerate}[label=(\alph*), nosep]
  \item $c_{\mathrm{low}}^{(0)} \le C^{*}_{\max}(\Gamma, m)$,
  \item $\mathrm{FFD}[\Gamma, C] \le m$ für jedes $C \ge c_{\mathrm{high}}^{(0)}$ und
  \item $C^{*}_{\max}(\Gamma, m) \le c_{\mathrm{high}}^{(0)}$.
\end{enumerate}
\end{lemma}

\begin{proof}
Aussage~(a) ist Lemma~\ref{lem:untere-schranken}.
Für~(b) sei angenommen, FFD eröffne bei Kapazität $C$ einen $(m+1)$-ten Bin, und zwar für die
Aufgabe $p_j$.
Zu diesem Zeitpunkt existieren genau $m$ Bins mit Lasten $l_1, \ldots, l_m$, und $p_j$ passt in
keinen davon, also $l_i + p_j > C$ für alle $i$ und damit $\sum_{i=1}^{m} l_i > m\,(C - p_j)$.
Da $p_j$ noch nicht platziert ist, gilt zugleich
$\sum_{i=1}^{m} l_i \le \sum_{t \ne j} p_t$.
Beides zusammen ergibt
\[
  C \;<\; \tfrac{1}{m}\textstyle\sum_{t} p_t + \tfrac{m-1}{m}\,p_j
  \;\le\; \tfrac{1}{m}\textstyle\sum_{t} p_t + p_{\max}
  \;=\; c_{\mathrm{high}}^{(0)} .
\]
Für $C \ge c_{\mathrm{high}}^{(0)}$ ist das ausgeschlossen, es entsteht also kein $(m+1)$-ter Bin.
Aussage~(c) folgt aus~(b): Bei Kapazität $c_{\mathrm{high}}^{(0)}$ erzeugt FFD höchstens $m$ Bins,
also ist $\mathrm{OPT}[\Gamma, c_{\mathrm{high}}^{(0)}] \le m$ und damit
$C^{*}_{\max}(\Gamma, m) \le c_{\mathrm{high}}^{(0)}$ nach Lemma~\ref{lem:dualitaet}.
\end{proof}

Die Güte von MULTIFIT wird über den folgenden Wert ausgedrückt, der die Tiefe der Binärsuche
ausklammert und allein das Verhältnis von FFD zum Optimum erfasst.

\begin{definition}[{Grenzverhältnis $r_m$, nach \cite[S.~79]{cao1995}}]
\label{def:rm}
Eine Instanz $\Gamma$ heißt $(p/q)$-\emph{Gegenbeispiel} für $m$ Maschinen, wenn
$\mathrm{FFD}[\Gamma, p] > m \ge \mathrm{OPT}[\Gamma, q]$ gilt.
Es sei
\[
  r_m \;=\; \sup\bigl\{\, p/q \;\bigm|\; \text{es gibt ein } (p/q)\text{-Gegenbeispiel für }
  m \text{ Maschinen} \,\bigr\}.
\]
\end{definition}

Kapitel~\ref{chap:performanz} beweist die Schranke $r_m \le 5/4$, ordnet die schärfere Schranke
$r_m \le 13/11$ in die Entwicklung der bekannten Ergebnisse ein und erläutert die Beweisidee von
\textcite{cao1995}; dieser Beweis selbst wird dort nicht in voller Länge geführt.
Ferner wird gezeigt, dass $13/11$ für $m \ge 13$ nicht verbessert werden kann.
Ebenfalls dort behandelt wird, wie aus $r_m$ die Güte $R_m(\mathrm{MF}(k)) \le r_m + 2^{-k}$
folgt; der additive Term geht auf die endliche Tiefe der Binärsuche zurück.

Für die Binärsuche ist die folgende Konsequenz von Definition~\ref{def:rm} entscheidend:
\begin{equation}
  C > r_m \cdot C^{*}_{\max}(\Gamma, m) \;\;\Longrightarrow\;\; \mathrm{FFD}[\Gamma, C] \le m.
  \label{eq:ffd-akzeptiert}
\end{equation}
Zum Beweis sei $q := C^{*}_{\max}(\Gamma, m)$ und $C > r_m \, q$.
Nach Lemma~\ref{lem:dualitaet} gilt $\mathrm{OPT}[\Gamma, q] \le m$, denn $q$ ist der optimale
Makespan.
Wäre nun $\mathrm{FFD}[\Gamma, C] > m$, so erfüllte $\Gamma$ beide Bedingungen aus
Definition~\ref{def:rm} und wäre ein $(C/q)$-Gegenbeispiel für $m$ Maschinen.
Da $r_m$ das Supremum aller Quotienten solcher Gegenbeispiele ist, folgte $C/q \le r_m$ -
im Widerspruch zu $C > r_m\, q$.
Also gilt $\mathrm{FFD}[\Gamma, C] \le m$.

MULTIFIT akzeptiert damit jede Kapazität oberhalb von $r_m \cdot C^{*}_{\max}(\Gamma, m)$ und kann
seine obere Intervallgrenze bis dorthin absenken.
Das ist die Eigenschaft, auf der die Schranke aus Kapitel~\ref{chap:performanz} beruht - und
die eine andere Subroutine an Stelle von FFD ebenfalls erfüllen muss, wenn dieselbe Schlussweise
greifen soll.

\begin{bemerkung}[Nicht-Monotonie von FFD]
\label{bem:nicht-monoton}
Die Binärsuche sucht keine Schwelle, denn $C \mapsto \mathrm{FFD}[\Gamma, C]$ fällt nicht monoton.
Für $\Gamma = (40, 39, 38, 36, 23, 22, 20, 19, 18, 15, 13, 13, 5)$ und $m = 4$ etwa gilt
$\mathrm{FFD}[\Gamma, 76] = 4$, aber $\mathrm{FFD}[\Gamma, 77] = 5$.
Bei $C = 76$ ist $39 + 38 = 77 > 76$; die Aufgabe~38 eröffnet daher einen dritten Bin, und es
entstehen die Lasten $(76, 75, 75, 75)$.
Bei $C = 77$ passen 39 und 38 gemeinsam in den zweiten Bin, der damit voll ist; am Ende bleibt eine
13er-Aufgabe allein übrig, und es entstehen die Lasten $(76, 77, 70, 65, 13)$.
Eine größere Kapazität kann also mehr Bins erfordern.
Weshalb die Binärsuche die Schranke aus Kapitel~\ref{chap:performanz} dennoch einhält, wird
dort ausgearbeitet.
\end{bemerkung}
